(1) Write $B = I \wedge u + C$ with $u \in I^\perp$ and $C \in \Lambda^2(I^\perp)$, which is possible and unique since $Q(I) \neq 0$ splits $V = \R I \oplus I^\perp$. Then, with bivectors acting on vectors as in \claimref{std-algebra.bivector-lie-algebra}, $B(I) = \ip{u}{I}I - Q(I)u + 0 = -Q(I)\,u$, which vanishes iff $u = 0$. So the stabiliser is $\Lambda^2(I^\perp)$, and $I^\perp$ carries the restricted form: by \claimref{synthesis.completion-recovers-the-real-clifford-algebras}(iii), signature $(1,4)$ if $I$ is a plus direction, $(2,3)$ if a minus direction. Over the scalars themselves, let $Q(I) = h^{2}$ with $h \neq 0$ and $\hat I = I/h$, so $Q(\hat I) = 1 = Q(e_{0})$ for the basis vector $e_{0}$ of square $1$. Since $Q(\hat I - e_{0}) + Q(\hat I + e_{0}) = 4$, one of $n = \hat I \mp e_{0}$ is non-null, and the reflection $x \mapsto -n x n^{-1}$ carries $\hat I$ to $\pm e_{0}$ and $I^\perp$ isometrically onto $e_{0}^\perp$, the standard space of signature $(1,4)$; conjugation by the versor $n$ carries $\Lambda^2(I^\perp)$ onto $\Lambda^2(e_{0}^\perp) = \mathfrak{so}(1,4)$. For $Q(I) = -h^{2}$ the same construction with a basis vector of square $-1$ gives $\mathfrak{so}(2,3)$. Triviality of the nilradical is the reductivity of orthogonal algebras of nondegenerate forms. (2) For $X \in M_I$: $\ip{Y}{I} = \ip{X}{I} - 1 = 0$ and $Q(Y) = Q(X) - 2\ip{X}{I}/Q(I) + 1/Q(I) = -1/Q(I)$. The map is an affine bijection commuting with the stabiliser. If $Q(I) > 0$: $I^\perp \cong \R^{1,4}$ and the locus $Q(Y) = -1/Q(I) < 0$ is the standard de Sitter hyperboloid, with induced signature $(1,3)$ and isometry group $O(1,4)$. If $Q(I) < 0$: $I^\perp \cong \R^{2,3}$ and the locus $Q(Y) > 0$ is the standard anti-de Sitter hyperboloid with isometry group $O(2,3)$. (3) Null rays of $I^\perp$: for $\R^{1,4}$ with form $t^2 - \sum_1^4 x_i^2$, the null cone is $t^2 = |x|^2$, whose rays are parametrised by the unit sphere $S^3 = \{|x| = 1, t = 1\}$; the induced degenerate form on the cone descends to a negative-definite form on the ray space, i.e.\ a spacelike (Riemannian) conformal $3$-sphere, the conformal infinity of de Sitter space. For $\R^{2,3}$ the same construction yields ray space $S^1 \times S^2$ (up to cover) carrying a Lorentzian conformal structure: the timelike infinity of anti-de Sitter space.
